\documentclass[12pt]{article}
\usepackage{amsmath, amssymb, amsthm}
\usepackage{enumerate}
\usepackage{geometry}
\geometry{margin=1in}

\newcommand{\Z}{\mathbb{Z}}

\begin{document}

\begin{center}
  {\large \textbf{MAT 211 --- Foundations of Mathematical Proof}}\\[4pt]
  {\large Final Exam}
\end{center}

\bigskip
\noindent\textbf{Name:} \hrulefill \hspace{2cm}
\textbf{Date:} \hrulefill

\bigskip
\begin{enumerate}

\item 
We define $\mathbb{Z}/26\mathbb{Z}$ to be the set of all congruence classes modulo $26$:
\[
\mathbb{Z}/26\mathbb{Z} = \{[0],[1],[2],\dots,[25]\},
\]
where $[a]$ denotes the set of all integers congruent to $a$ modulo $26$.

\begin{enumerate}
    \item Find the additive inverse of $[7]$ in $\Z/26\Z$ and give a general
            formula for the additive inverse of $[a]$ in $\Z/26\Z$.    

    \item Does $[4]$ have a multiplicative inverse in $\Z/26\Z$?
    Explain briefly why or why not.
    
    \item Find the multiplicative inverse of $[3]$ in $\Z/26\Z$ using
    the extended Euclidean algorithm or Bézout's Identity.
    Verify your answer.

\end{enumerate}

\bigskip

\item \textbf{Negating Quantifiers and Algebraic Structure}

\begin{enumerate}[(a)]

\item Let $R$ be an algebraic ring.  Consider the statement:
\[
(\forall a \in R,\; a \neq 0 \Rightarrow \exists b \in R \text{ such that } ab = 1).
\]
Write the negation of this statement symbolically using quantifiers.

\item Interpret both the original statement and its negation in words.
What algebraic property does the original statement describe?

\item Consider $R = \mathbb{Z}/n\mathbb{Z}$.
\begin{itemize}
\item Give an example of $n$ for which the original statement is true.
\item Give an example of $n$ for which the negation is true.
\end{itemize}
Briefly justify your answers.

\end{enumerate}

\bigskip

\item \textbf{Breaking an Affine Cipher}

\textit{Imagine you are participating in a cyber-security hack-a-thon competition.
Suppose you find an intercepted message encrypted using an \textbf{affine cipher}}
\[
  f(x) \equiv ax + b \pmod{26},
\]
\textit{where letters are encoded as integers $0$ through $25$
(\texttt{A} $=0$, \ldots, \texttt{Z} $=25$), and $a,b \in \mathbb{Z}/26\mathbb{Z}$.}

\medskip
\textit{Your team identifies the following plaintext--ciphertext pairs:}
\[
\texttt{H} \mapsto \texttt{M}, \qquad \texttt{E} \mapsto \texttt{X}.
\]

\begin{enumerate}[(a)]

\item Determine the values of $a$ and $b$ in $\mathbb{Z}/26\mathbb{Z}$ used in the encryption.
Show your work clearly.

\item Explain why solving for $a$ requires that a certain integer be
invertible modulo $26$.  Identify this integer and justify why it
\emph{does} have a multiplicative inverse.

\item State the condition on $a$ that guarantees the cipher is
decryptable (i.e., that $f$ has an inverse on $\mathbb{Z}/26\mathbb{Z}$),
and verify that your value of $a$ satisfies this condition.

\item Your team now uses the hacked encryption rule to analyze a second
intercepted message.  Decrypt
\[
\texttt{IVVL}.
\]

\bigskip 

\end{enumerate}

\item \textbf{Multiplicative Inverses and Fermat's Little Theorem}

\textit{Fermat's Little Theorem states: if $p$ is prime and
$\gcd(a,p)=1$, then $a^{p-1} \equiv 1 \pmod{p}$.}

\begin{enumerate}[(a)]

\item Use Fermat's Little Theorem to show that if $p$ is prime and
$\gcd(a,p)=1$, then $a$ has a multiplicative inverse modulo $p$.
Identify an explicit expression for $a^{-1} \pmod{p}$.

\item Use your result to compute $3^{-1} \pmod{7}$.

\item Explain why this method cannot be used to find multiplicative
inverses modulo $26$.  What property of $26$ is responsible, and why
does it matter?

\end{enumerate}

\bigskip 

\item \textbf{The Discrete Logarithm Problem}

\textit{You encounter another cryptographic challenge.
This time the system is based on \textbf{modular exponentiation}.
A server publishes a prime $p = 29$ and a base $g = 2$.
One party chooses a secret exponent $x$ and transmits only}
\[
  h \equiv g^x \equiv 2^x \pmod{29}.
\]
\textit{Your team intercepts $h = 7$.  To break the system you must
recover $x$.}

\begin{enumerate}[(a)]

\item By computing successive powers $2^1, 2^2, 2^3, \dots \pmod{29}$,
find the smallest positive integer $x$ such that
$2^x \equiv 7 \pmod{29}$.  (Organizing your work in a table may help.)

\item Fermat's Little Theorem states that
$2^{28} \equiv 1 \pmod{29}$.
Explain why this implies
\[
2^{x+28} \equiv 2^x \pmod{29},
\]
and therefore why it is sufficient to search only
$x \in \{1,2,\dots,28\}$.

\item Suppose instead that $p$ is a prime with about $200$ digits.
Roughly how many possible values of $x$ would need to be checked in a
brute-force search?  Explain why such a search is computationally
infeasible, even for modern computers.

\newpage 


\item Contrast the two directions:
\begin{itemize}
\item \textbf{Forward:} Given $g$, $x$, and $p$, compute $g^x \pmod{p}$.
\item \textbf{Backward:} Given $g$, $h$, and $p$, find $x$ such that
$g^x \equiv h \pmod{p}$.
\end{itemize}
Explain intuitively why the forward direction is easier than the
backward direction.  Why is this kind of asymmetry (often called a
\emph{one-way function}) useful in cryptography?

\end{enumerate}

\bigskip

\item \textbf{A Function from \(\mathbb{Z}^2\) to \(\mathbb{Z}\)}

Let \(a,b \in \mathbb{Z}\), not both zero, and define
\[
f : \mathbb{Z}^2 \to \mathbb{Z}, \qquad f(m,n)=am+bn.
\]

\begin{enumerate}[(a)]

\item Describe the range of \(f\).  Your answer should be written in terms of \(\gcd(a,b)\).

\item Determine precisely when \(f\) is surjective.

\item Determine whether \(f\) can ever be injective.  Justify your answer.

\end{enumerate}

\bigskip 

\item \textbf{Prove:} Let $a,b \in \mathbb{Z}$. Show that if a prime $p$ divides $ab$, then $p$ divides $a$ or $p$ divides $b$.

\bigskip 

\item \textbf{Prove:} Let $a,b \in \mathbb{Z}$. Prove that if $ab$ is odd, then $a^2+b^2$ is even. 

\bigskip

\item \textbf{Prove:} Let $a,b \in \mathbb{R}$.  Prove that if $a \leq b$, then $a^3 \leq b^3$.

\bigskip 

\item \textbf{Reflection:}  The proof for question \#9 (above) involves a statement which was also proven using the programming language Lean in the LEAN4 homework (where it was also question \#9). 
\begin{enumerate}
    \item Describe any differences between how you proved the statement in Lean (in the previous assignment) and how you proved the statement above, on this assignment. How was your proof strategy different?
    \item What benefits or drawbacks do you find when comparing traditional proof writing with writing proofs in Lean?
\end{enumerate}  
\end{enumerate}

\end{document}